% report/ablation_section.tex -- draft ablation section for merging into the main report (IEEEtran). Not included anywhere yet.
% Requires \usepackage{booktabs}. Tables are copied verbatim from results/ablation/tables.tex (same style: caption above, \footnotesize, booktabs).
% Every number in this file is checked against results/ablation/*.json by ablation/verify_findings.py (one sentence per line on purpose).

\subsection{From the original notebook to v2}
\label{sec:ablation-ladder}

We rebuilt the path from the original notebook to the final pipeline as a ladder of single changes and scored every rung on the same corrected test sets.
Task~2 rungs are trained with three seeds (42, 43, 44) and reported as mean $\pm$ sample standard deviation; Task~1 and Task~3 involve no training and are single deterministic runs.
Across the Task~2 rungs the seed standard deviation of test macro-F1 lies between 0.001 and 0.012, and the standard deviation of paired same-seed differences between 0.002 and 0.019, so we treat differences below about 0.01 as not established.

\paragraph{Task 2.}
Table~\ref{tab:t2_ladder} lists the rungs; all are evaluated on the full corrected test set of 2237 images.
The path fix is the only established accuracy gain.
The original notebook's \texttt{lstrip('data\_image/')} strips a set of characters instead of a prefix and silently dropped both earthquake events from the training and development data (L0).
Restoring them raises test macro-F1 from 0.494 $\pm$ 0.006 to 0.588 $\pm$ 0.012 (L1), a paired gain of $+0.095 \pm 0.011$ that is positive at all three seeds.
Every later step moves test macro-F1 by less than the noise floor: the paired differences are $+0.004 \pm 0.002$ (QuickGELU CLIP, L2), $+0.008 \pm 0.007$ (10-epoch cap, L3), $-0.013 \pm 0.015$ (deterministic cached features, L3b), $+0.012 \pm 0.019$ (mixed precision, L4a) and $-0.011 \pm 0.010$ (randomly initialised text encoder, L4).
The rest of v2 therefore delivers equal accuracy at a fraction of the compute: L1 reaches 0.588 $\pm$ 0.012 and v2 (L4) 0.587 $\pm$ 0.001 test macro-F1, at 9.8 instead of 66.4 minutes of wall-clock time per run, about one seventh, measured on a single laptop GPU.
The whole bundle from L3 to L4 is consistently slightly lower ($-0.013 \pm 0.005$, negative at every seed), but no single step of it is established, including the effect of the random initialisation, which is no larger than the opposite-signed effect of mixed precision alone.
Two implementation facts explain details of the ladder.
The original notebook applied a random crop to development and test images as well as to training images, which is why its CLIP features could not be cached before L3b.
v2's text encoder starts from random weights instead of pretrained DistilBERT: the cosine similarity between its saved word embeddings and the pretrained ones is $-0.0002$.

\begin{table}[t]
\centering
\footnotesize
\caption{Task 2 ladder from the original notebook (L0) to v2 (L4). All rungs are scored on the full corrected test set (2237 rows).}
\label{tab:t2_ladder}
\begin{tabular}{llcllll}
\toprule
Rung & Change & $n$ & Dev macro-F1 & Acc. & Macro-F1 & Weighted-F1 \\
\midrule
L0 & v1 buggy train/dev, GELU CLIP, random crop, 30 ep & 3 & 0.557 $\pm$ 0.016 & 0.767 $\pm$ 0.002 & 0.494 $\pm$ 0.006 & 0.760 $\pm$ 0.002 \\
L1 & + full corrected train/dev & 3 & 0.643 $\pm$ 0.044 & 0.788 $\pm$ 0.008 & 0.588 $\pm$ 0.012 & 0.788 $\pm$ 0.011 \\
L2 & + QuickGELU CLIP & 3 & 0.648 $\pm$ 0.043 & 0.786 $\pm$ 0.012 & 0.592 $\pm$ 0.012 & 0.787 $\pm$ 0.013 \\
L3 & + epoch budget 30 $\to$ 10 & 3 & 0.595 $\pm$ 0.008 & 0.797 $\pm$ 0.005 & 0.600 $\pm$ 0.005 & 0.799 $\pm$ 0.004 \\
L3b & + deterministic transform, cached features & 3 & 0.581 $\pm$ 0.011 & 0.792 $\pm$ 0.008 & 0.587 $\pm$ 0.011 & 0.793 $\pm$ 0.007 \\
L4a & + AMP (pretrained DistilBERT) & 3 & 0.585 $\pm$ 0.008 & 0.799 $\pm$ 0.006 & 0.598 $\pm$ 0.010 & 0.800 $\pm$ 0.006 \\
L4 & v2: random-init DistilBERT & 3 & 0.581 $\pm$ 0.003 & 0.793 $\pm$ 0.005 & 0.587 $\pm$ 0.001 & 0.795 $\pm$ 0.002 \\
\bottomrule
\end{tabular}
\end{table}

\paragraph{Task 3.}
Table~\ref{tab:t3_ladder} follows the nearest-neighbour index from the original notebook to v2; every rung is scored on the same 529 test queries and the three real damage classes.
The original index held all 18082 annotated images, including the test images, and 14556 of them (80.5\%) carried a ``Don't know'' label that never occurs in the test split.
The original notebook's score is not its real performance: with the test images in the index, all 529 test queries retrieve themselves among their five nearest neighbours, and the resulting macro-F1 of 0.688 (accuracy 0.722) falls to 0.378 (accuracy 0.503) once they are removed (L1).
L1 is therefore the fair measure of the original system; removing the polluted labels recovers 0.504 (L2), and the last two steps, dropping development images (L3) and same-tweet siblings (L4, v2), give 0.509 and 0.491, which are within the noise of 529 queries.

\begin{table}[t]
\centering
\footnotesize
\caption{Task 3 ladder (deterministic, $n=1$ per rung). Every rung is scored on the same 529 test queries (3 real classes).}
\label{tab:t3_ladder}
\begin{tabular}{llrrrr}
\toprule
Rung & Index & Rows & Acc. & Macro-F1 & Weighted-F1 \\
\midrule
L0 & v1 index (incl. test, ``Don't know'' pollution) & 18082 & 0.7221 & 0.6880 & 0.7630 \\
L1 & test rows removed & 17553 & 0.5028 & 0.3777 & 0.5328 \\
L2 & pollution removed (train+dev) & 2997 & 0.6767 & 0.5037 & 0.6318 \\
L3 & dev removed (train only) & 2468 & 0.6749 & 0.5091 & 0.6309 \\
L4 & tweet siblings removed (= v2) & 2361 & 0.6673 & 0.4914 & 0.6202 \\
\bottomrule
\end{tabular}
\end{table}

\subsection{Component ablation}
\label{sec:ablation-components}

\paragraph{Task 2.}
Table~\ref{tab:t2_partb} removes one component at a time from the L4a configuration (three seeds each).
A 200k-parameter MLP on the cached CLIP vector scores 0.601 $\pm$ 0.005 test macro-F1 against 0.598 $\pm$ 0.010 for the full 67M-parameter model, a paired difference of $+0.003 \pm 0.009$, and concatenating the two streams without attention gives $-0.003 \pm 0.012$.
Neither difference is established.
The text stream alone reaches only 0.294 $\pm$ 0.014 ($-0.305 \pm 0.012$ against the full model), which is established.
For this task the text is redundant with the image: the target is the image's own humanitarian label, and the tweet text adds nothing measurable once the CLIP vector is available.
Removing the text-to-image attention block gives $-0.029 \pm 0.024$ (negative at all three seeds), which we regard as suggestive only.
That block attends over a single image vector, so its attention weight is identically one and it contributes an additional linear copy of the image; the image-to-text block and the direct residual image path remain in the ablated model.

\begin{table}[t]
\centering
\footnotesize
\caption{Task 2 component ablation on L4a's configuration (10 epochs, cached CLIP, AMP, pretrained DistilBERT).}
\label{tab:t2_partb}
\begin{tabular}{lcrllll}
\toprule
Variant & $n$ & Params & Dev macro-F1 & Acc. & Macro-F1 & Weighted-F1 \\
\midrule
full model (L4a) & 3 & 67\,352\,840 & 0.585 $\pm$ 0.008 & 0.799 $\pm$ 0.006 & 0.598 $\pm$ 0.010 & 0.800 $\pm$ 0.006 \\
text only & 3 & 66\,628\,104 & 0.295 $\pm$ 0.005 & 0.486 $\pm$ 0.010 & 0.294 $\pm$ 0.014 & 0.499 $\pm$ 0.008 \\
image only & 3 & 199\,688 & 0.568 $\pm$ 0.008 & 0.788 $\pm$ 0.002 & 0.601 $\pm$ 0.005 & 0.793 $\pm$ 0.002 \\
concat, no attention & 3 & 66\,825\,480 & 0.580 $\pm$ 0.012 & 0.796 $\pm$ 0.007 & 0.595 $\pm$ 0.009 & 0.795 $\pm$ 0.007 \\
no text$\to$image attention & 3 & 67\,089\,160 & 0.583 $\pm$ 0.015 & 0.795 $\pm$ 0.005 & 0.570 $\pm$ 0.014 & 0.795 $\pm$ 0.005 \\
\bottomrule
\end{tabular}
\end{table}

\paragraph{Task 3.}
Tables~\ref{tab:t3_k}--\ref{tab:t3_emb} vary one factor of the nearest-neighbour classifier at a time on the tweet-disjoint index (single deterministic runs on 529 test queries).
The value $k=5$ is the best on the development set for the headline configuration (macro-F1 0.483), and test macro-F1 lies between 0.454 and 0.515 across $k \in \{1, 3, 5, 10, 20\}$, which is not an established effect.
Similarity-weighted voting gives metrics identical to majority voting in all 15 combinations of embedding and $k$, because the neighbours of a query have similar similarities (mean gap between the first and fifth neighbour 0.046).
Class-balanced voting raises development macro-F1 by 0.037 but test macro-F1 by only 0.004, and lowers accuracy from 0.667 to 0.586.
Image-only retrieval (0.507) is not worse than the fused embedding (0.491), and text-only retrieval is clearly weaker (0.428).
A logistic-regression probe on the same embeddings scores higher than five-nearest-neighbour voting on the fused embedding (0.615 against 0.491) and on the image embedding (0.595 against 0.507); this is large relative to the noise floor but it is a single run and has not been replicated across seeds.

\begin{table}[t]
\centering
\footnotesize
\caption{Task 3: effect of $k$ (fused embedding, majority vote, tweet-disjoint index). Headline is $k=5$. Deterministic, $n=1$.}
\label{tab:t3_k}
\begin{tabular}{rrrr}
\toprule
$k$ & Dev macro-F1 & Test acc. & Test macro-F1 \\
\midrule
1 & 0.4627 & 0.6011 & 0.4538 \\
3 & 0.4739 & 0.6749 & 0.5150 \\
5 & 0.4832 & 0.6673 & 0.4914 \\
10 & 0.4438 & 0.6843 & 0.5006 \\
20 & 0.4334 & 0.6843 & 0.4704 \\
\bottomrule
\end{tabular}
\end{table}


\begin{table}[t]
\centering
\footnotesize
\caption{Task 3: voting rule ($k=5$, fused embedding). Deterministic, $n=1$.}
\label{tab:t3_vote}
\begin{tabular}{lrrr}
\toprule
Vote & Dev macro-F1 & Test acc. & Test macro-F1 \\
\midrule
majority (v2) & 0.4832 & 0.6673 & 0.4914 \\
similarity-weighted & 0.4832 & 0.6673 & 0.4914 \\
class-balanced & 0.5198 & 0.5860 & 0.4954 \\
\bottomrule
\end{tabular}
\end{table}


\begin{table}[t]
\centering
\footnotesize
\caption{Task 3: embedding and classifier (index/training rows identical; logistic-regression $C$ tuned on dev). Deterministic, $n=1$.}
\label{tab:t3_emb}
\begin{tabular}{llrrr}
\toprule
Embedding & Classifier & Dev macro-F1 & Test acc. & Test macro-F1 \\
\midrule
image & 5-NN majority & 0.5427 & 0.6749 & 0.5074 \\
image & logistic regression & 0.6208 & 0.6730 & 0.5946 \\
text & 5-NN majority & 0.4021 & 0.6068 & 0.4276 \\
text & logistic regression & 0.4731 & 0.5085 & 0.4632 \\
fused & 5-NN majority & 0.4832 & 0.6673 & 0.4914 \\
fused & logistic regression & 0.5926 & 0.6843 & 0.6146 \\
\bottomrule
\end{tabular}
\end{table}

\subsection{Robustness}
\label{sec:ablation-robustness}

\paragraph{Missing modality.}
Table~\ref{tab:c1} replaces one input at test time without retraining: a grey image, or an empty string.
A grey image collapses both Task~2 models to 0.080 (L4) and 0.081 (L4a) macro-F1, within 0.002 of the score of always predicting the majority class (0.0801) at every seed, so the image carries essentially all of the decision.
Removing the text has a model-dependent effect: v2's model (L4) is unchanged ($+0.003$, 0.587 to 0.590), whereas the model with pretrained DistilBERT (L4a) loses between 0.058 and 0.103 macro-F1 across the three seeds (mean $-0.087$).
We read this as redundancy and not as an ignored input: text-only training reaches 0.294 and image-only training matches the full model, so the text does not supply information that the image lacks.
An empty string is also an unusual input for DistilBERT: training texts are never empty, so the encoder sees only the start and end tokens, and the L4a drop may reflect that input shift as much as lost evidence.
For Task~3 a grey image costs 0.039 macro-F1 and an empty text 0.052 (single run).

\begin{table}[t]
\centering
\footnotesize
\caption{Missing-modality robustness at test time (no retraining). $\Delta$ is the macro-F1 change from the full-input score.}
\label{tab:c1}
\begin{tabular}{llcrrr}
\toprule
Model & Condition & $n$ & Acc. & Macro-F1 & $\Delta$ macro-F1 \\
\midrule
T2 L4a & full input & 3 & 0.799 $\pm$ 0.006 & 0.598 $\pm$ 0.010 & -- \\
 & grey image & 3 & 0.472 $\pm$ 0.001 & 0.081 $\pm$ 0.001 & -0.518 \\
 & empty text & 3 & 0.774 $\pm$ 0.005 & 0.511 $\pm$ 0.034 & -0.087 \\
T2 L4 (v2) & full input & 3 & 0.793 $\pm$ 0.005 & 0.587 $\pm$ 0.001 & -- \\
 & grey image & 3 & 0.472 $\pm$ 0.000 & 0.080 $\pm$ 0.000 & -0.507 \\
 & empty text & 3 & 0.796 $\pm$ 0.003 & 0.590 $\pm$ 0.013 & +0.003 \\
T3 5-NN & full input & 1 & 0.667 & 0.491 & -- \\
 & grey image & 1 & 0.565 & 0.452 & -0.039 \\
 & empty text & 1 & 0.665 & 0.440 & -0.052 \\
\bottomrule
\end{tabular}
\end{table}

\paragraph{Leave-one-disaster-out.}
Tables~\ref{tab:c2_t1} and~\ref{tab:c2_t3} train (Task~1) or build the index (Task~3) on six events and test on the seventh; the held-out event contributes no rows, which the code asserts for all 14 folds.
Averaged over the seven events, Task~1 macro-F1 falls from 0.714 (all events seen) to 0.680 and Task~3 from 0.381 to 0.354.
The Task~1 drop appears in 6 of 7 events (largest: Sri Lanka floods, 0.932 to 0.793) and the Task~3 drop in 5 of 7, but three Task~3 events have only 13 to 23 test images, so we regard the Task~3 figures as indicative only.
Task~2 was not evaluated this way and is left as future work.

\begin{table}[t]
\centering
\footnotesize
\caption{Leave-one-disaster-out, T1 (train/index on the other six events' train split; deterministic, $n=1$ per fold). Per-event macro-F1 averages the classes present in that event. In-dist.\ = standard model on the same rows.}
\label{tab:c2_t1}
\begin{tabular}{lrrrr}
\toprule
Held-out event & Test rows & Acc. & LODO macro-F1 & In-dist. macro-F1 \\
\midrule
california\_wildfires & 240 & 0.708 & 0.615 & 0.606 \\
hurricane\_harvey & 568 & 0.805 & 0.735 & 0.758 \\
hurricane\_irma & 548 & 0.776 & 0.675 & 0.697 \\
hurricane\_maria & 542 & 0.723 & 0.692 & 0.742 \\
iraq\_iran\_earthquake & 73 & 0.753 & 0.622 & 0.625 \\
mexico\_earthquake & 159 & 0.711 & 0.628 & 0.638 \\
srilanka\_floods & 107 & 0.804 & 0.793 & 0.932 \\
\textbf{macro average} & 2237 & 0.754 & 0.680 & 0.714 \\
\bottomrule
\end{tabular}
\end{table}


\begin{table}[t]
\centering
\footnotesize
\caption{Leave-one-disaster-out, T3 (train/index on the other six events' train split; deterministic, $n=1$ per fold). Per-event macro-F1 averages the classes present in that event. In-dist.\ = standard model on the same rows.}
\label{tab:c2_t3}
\begin{tabular}{lrrrr}
\toprule
Held-out event & Test rows & Acc. & LODO macro-F1 & In-dist. macro-F1 \\
\midrule
california\_wildfires & 87 & 0.862 & 0.309 & 0.315 \\
hurricane\_harvey & 132 & 0.629 & 0.481 & 0.494 \\
hurricane\_irma & 120 & 0.408 & 0.271 & 0.483 \\
hurricane\_maria & 133 & 0.579 & 0.389 & 0.343 \\
iraq\_iran\_earthquake & 23 & 0.783 & 0.293 & 0.302 \\
mexico\_earthquake & 21 & 0.857 & 0.462 & 0.417 \\
srilanka\_floods & 13 & 0.692 & 0.273 & 0.316 \\
\textbf{macro average} & 529 & 0.687 & 0.354 & 0.381 \\
\bottomrule
\end{tabular}
\end{table}

\paragraph{What is established.}
With three seeds and a noise floor of about 0.01 macro-F1, we regard only the following as established: the path-fix gain ($+0.095$), the collapse of the text-only model, the collapse of every Task~2 model without its image, and the fall of the original Task~3 score once self-retrieval is removed.
Differences of about 0.01 between configurations are not established; this includes the effect of the random initialisation in v2 and of mixed precision, the comparison between image-only, concatenation and the full model, and all differences among the Task~3 rungs from L2 to L4.
